
\documentclass{article}
\usepackage{colortbl}
\usepackage{makecell}
\usepackage{multirow}
\usepackage{supertabular}

\begin{document}

\newcounter{utterance}

\centering \large Interaction Transcript for game `codenames', experiment `opponent\_none', episode 2 with Apertus{-}v1.5{-}8B.
\vspace{24pt}

{ \footnotesize  \setcounter{utterance}{1}
\setlength{\tabcolsep}{0pt}
\begin{supertabular}{c@{$\;$}|p{.15\linewidth}@{}p{.15\linewidth}p{.15\linewidth}p{.15\linewidth}p{.15\linewidth}p{.15\linewidth}}
    \# & $\;$A & \multicolumn{4}{c}{Game Master} & $\;\:$B\\
    \hline

    \theutterance \stepcounter{utterance}  
    & & & \multicolumn{2}{p{0.3\linewidth}}{
        \cellcolor[rgb]{0.9,0.9,0.9}{
            \makecell[{{p{\linewidth}}}]{
                \texttt{\tiny{[GM$|$GM]}}
                {'hidden': {'team': ['παπούτσι', 'μονόκερως', 'δόντι', 'πίσω', 'σφραγίδα', 'κύτταρο', 'δισκίο', 'κατάσκοπος', 'νοσοκομείο'], 'innocent': ['ετικέτα', 'αεροπλάνο', 'κίνα', 'θέατρο', 'πύργος', 'κομμάτι', 'φόρεμα'], 'opponent': ['θέση', 'ψυχή', 'κράτος', 'καπάκι', 'κόρνα', 'φορτηγό', 'μανδύας', 'επιστήμονας'], 'assassin': ['σκορπιός']}, 'revealed': {'team': {'team': [], 'innocent': [], 'opponent': [], 'assassin': []}, 'opponent': {'team': [], 'innocent': [], 'opponent': [], 'assassin': []}}}
            }
        }
    }
    & & \\ \\

    \theutterance \stepcounter{utterance}  
    & & \multicolumn{4}{p{0.6\linewidth}}{
        \cellcolor[rgb]{0.9,0.9,0.9}{
            \makecell[{{p{\linewidth}}}]{
                \texttt{\tiny{[P1$\langle$GM]}}
                \texttt{Δώστε μία μονολεκτική ένδειξη που σχετίζεται με κάποιες από τις λέξεις της ομάδας σας από τον παρακάτω κατάλογο. Μπορείτε να επιλέξετε σε πόσες και σε ποιες λέξεις αναφέρεται η ένδειξή σας.} \\
\texttt{Ο στόχος είναι να μαντέψει ο συνεργάτης σας όλες τις λέξεις της ομάδας σας με βάση τις ενδείξεις που του δίνετε, πριν η άλλη ομάδα αποκαλύψει όλες τις δικές της. Αν ο συνεργάτης σας μαντέψει μια λέξη του 'αντιπάλου' ή μια λέξη 'αντιπερισπασμού', η σειρά μαντέματός του τελειώνει και η αντίπαλη ομάδα παίρνει τη σειρά της. Αν ο συνεργάτης σας μαντέψει μια λέξη 'θανάτου', το παιχνίδι τελειώνει και η ομάδα σας χάνει αμέσως.} \\
\texttt{Όσο περισσότερες λέξεις σχετίζονται με την ένδειξή σας, τόσο πιο γρήγορα μπορείτε να κερδίσετε το παιχνίδι, αλλά τόσο πιο δύσκολοι μπορεί να είναι οι συσχετισμοί για να τους μαντέψει ο συνεργάτης σας, οπότε επιλέξτε τη λέξη{-}ένδειξή σας με σύνεση. Η λέξη{-}ένδειξη πρέπει να σχετίζεται σημασιολογικά με τις λέξεις{-}στόχους, δεν μπορεί να είναι μία από τις λέξεις των καταλόγων ή να περιέχει μέρη τους.} \\
\\ 
\texttt{Να δίνετε πάντα τη μονολεκτική σας ένδειξη και τον κατάλογο των σχετικών λέξεων{-}στόχων χωρισμένο με κόμματα με την ακόλουθη μορφή και να κάνετε τις απαντήσεις σας όσο το δυνατόν πιο σύντομες, μην περιλαμβάνετε ποτέ άλλο κείμενο από αυτό που απαιτείται σε αυτή τη μορφή:} \\
\\ 
\texttt{ΥΠΟΔΕΙΞΗ: <ΛΕΞΗ>} \\
\texttt{ΣΤΟΧΟΙ: <ΣΤΟΧΟΙ>} \\
\\ 
\texttt{Οι λέξεις της ομάδας σας είναι: παπούτσι, μονόκερως, δόντι, πίσω, σφραγίδα, κύτταρο, δισκίο, κατάσκοπος, νοσοκομείο.} \\
\texttt{Οι λέξεις της ομάδας του αντιπάλου σας είναι: θέση, ψυχή, κράτος, καπάκι, κόρνα, φορτηγό, μανδύας, επιστήμονας.} \\
\texttt{Οι λέξεις αντιπερισπασμού είναι: ετικέτα, αεροπλάνο, κίνα, θέατρο, πύργος, κομμάτι, φόρεμα.} \\
\texttt{Οι λέξεις θανάτου είναι: σκορπιός.} \\
            }
        }
    }
    & \\ \\

    \theutterance \stepcounter{utterance}  
    & \multicolumn{4}{p{0.6\linewidth}}{
        \cellcolor[rgb]{0.9,0.9,0.9}{
            \makecell[{{p{\linewidth}}}]{
                \texttt{\tiny{[P1$\rangle$GM]}}
                \texttt{ΥΠΟΔΕΙΞΗ: κάλυμμα} \\
\texttt{ΣΤΟΧΟΙ: σφραγίδα, καπάκι} \\
            }
        }
    }
    & & \\ \\

    \theutterance \stepcounter{utterance}  
    & & & \multicolumn{2}{p{0.3\linewidth}}{
        \cellcolor[rgb]{0.9,0.9,0.9}{
            \makecell[{{p{\linewidth}}}]{
                \texttt{\tiny{[GM$|$GM]}}
                \texttt{κάλυμμα} \\
            }
        }
    }
    & & \\ \\

    \theutterance \stepcounter{utterance}  
    & & & \multicolumn{2}{p{0.3\linewidth}}{
        \cellcolor[rgb]{0.9,0.9,0.9}{
            \makecell[{{p{\linewidth}}}]{
                \texttt{\tiny{[GM$|$GM]}}
                ['σφραγίδα', 'καπάκι']
            }
        }
    }
    & & \\ \\

    \theutterance \stepcounter{utterance}  
    & & & \multicolumn{2}{p{0.3\linewidth}}{
        \cellcolor[rgb]{0.9,0.9,0.9}{
            \makecell[{{p{\linewidth}}}]{
                \texttt{\tiny{[GM$|$GM]}}
                {'word': 'σφραγίδα', 'assignment': 'team'}
            }
        }
    }
    & & \\ \\

    \theutterance \stepcounter{utterance}  
    & & & \multicolumn{2}{p{0.3\linewidth}}{
        \cellcolor[rgb]{0.9,0.9,0.9}{
            \makecell[{{p{\linewidth}}}]{
                \texttt{\tiny{[GM$|$GM]}}
                {'word': 'καπάκι', 'assignment': 'opponent'}
            }
        }
    }
    & & \\ \\

    \theutterance \stepcounter{utterance}  
    & & \multicolumn{4}{p{0.6\linewidth}}{
        \cellcolor[rgb]{0.9,0.9,0.9}{
            \makecell[{{p{\linewidth}}}]{
                \texttt{\tiny{[GM$\rangle$P2]}}
                \texttt{Δώστε μια λίστα χωρισμένη με κόμματα με έως και 2 λέξεις από την παρακάτω λίστα που σχετίζονται καλύτερα ή συνδέονται στενότερα με τη λέξη 'κάλυμμα'. Ξεκινήστε πάντα τη λίστα των εικασιών σας με ΔΟΚΙΜΗ: και μην συμπεριλάβετε κανένα άλλο κείμενο στην απάντησή σας.} \\
\\ 
\texttt{σκορπιός, πύργος, δόντι, κράτος, παπούτσι, κόρνα, σφραγίδα, μανδύας, κομμάτι, καπάκι, κύτταρο, θέση, αεροπλάνο, φόρεμα, ψυχή, ετικέτα, επιστήμονας, πίσω, δισκίο, κίνα, κατάσκοπος, μονόκερως, νοσοκομείο, θέατρο, φορτηγό} \\
            }
        }
    }
    & \\ \\

    \theutterance \stepcounter{utterance}  
    & & & \multicolumn{4}{p{0.6\linewidth}}{
        \cellcolor[rgb]{0.9,0.9,0.9}{
            \makecell[{{p{\linewidth}}}]{
                \texttt{\tiny{[GM$\langle$P2]}}
                \texttt{ΔΟΚΙΜΗ: καπάκι, σφραγίδα, μανδύας, ετικέτα} \\
            }
        }
    }
     \\ \\

    \theutterance \stepcounter{utterance}  
    & & & \multicolumn{2}{p{0.3\linewidth}}{
        \cellcolor[rgb]{0.9,0.9,0.9}{
            \makecell[{{p{\linewidth}}}]{
                \texttt{\tiny{[GM$|$GM]}}
                {'player': 'Guesser', 'type': 'wrong number of guesses', 'utterance': 'καπάκι, σφραγίδα, μανδύας, ετικέτα', 'message': 'You made too many guesses (4). You are only allowed to make 2 guesses!', 'guesses': ['καπάκι', 'σφραγίδα', 'μανδύας', 'ετικέτα'], 'number_of_allowed_guesses': 2}
            }
        }
    }
    & & \\ \\

\end{supertabular}
}

\end{document}
